\documentclass{article}
% This is the research proposal for Groningen University. The submission of RP is the necessary 
% step before master's thesis registration.
% The proposal should not be longer than 10 pages (before methodology section, it is already 6 pages long)
\usepackage[margin=1in]{geometry}

% Math and figures
\usepackage{amsmath,amssymb}
\usepackage{graphicx}
\usepackage{tikz-cd}
\usepackage{multicol}

\usepackage{microtype}
\usepackage{listings}
\usepackage{xcolor}
\usepackage{booktabs}
\usepackage{hyperref}
\usepackage{enumitem}
\usepackage{mdframed}

% Paragraphs
\setlength{\parindent}{0pt}
\setlength{\parskip}{1\baselineskip}

\usepackage{natbib}
\usepackage{url}
\usepackage{listings}

\bibliographystyle{plainnat}

\title{Representing Self-Repair In Meaning Representation System: A Case Study of SBN}
\author{Hongxu Zhou}


\begin{document}
\pagestyle{empty}

\maketitle

\hypersetup{
    colorlinks=true,
    linkcolor=black,
    citecolor=black,
    urlcolor=blue
}

% Style for SBN sequences
% Save this for the thesis
% --Hongxu Zhou
\lstdefinestyle{sbn}{
    basicstyle=\ttfamily\small,
    breaklines=false,
    frame=single,
    rulecolor=\color{gray!60},
    backgroundcolor=\color{gray!6},
    xleftmargin=1.5em,
    xrightmargin=1.5em,
    framexleftmargin=1em,
    columns=flexible,
    keepspaces=true,
    showstringspaces=false,
    aboveskip=0.6em,
    belowskip=0.6em,
    commentstyle=\color{gray!80}\itshape,
    morecomment=[l]{\%},
}

\newcommand{\sbnterm}[1]{\texttt{#1}}



\section{Introduction}

Self-repair is a routine property of spontaneous speech, where a speaker breaks off
an utterance and revises part of what they have just said. Computational work
generally follows the taxonomy of \citet{shribergPreliminaries1994}, which
%% TODO: add Shriberg 1994, Levelt 1983 and Fox, Maschler & Uhmann 2010
%% (J. Pragmatics 42(9), key foxCrosslinguisticStudySelfrepair2010) to .bib
builds on \citet{leveltMonitoringSelfrepair1983} and decomposes a repair into
three parts: the \emph{reparandum}, the material to
be revised; an \emph{interregnum}, an editing phrase such as \textit{I mean} that
may or may not be present; and the \emph{repair}, the material that replaces the reparandum.
What the reparandum contributes differs across the three repair types the
taxonomy distinguishes. A repetition restates it and a restart abandons it, so in
neither case does removing it change what the speaker has committed to.
Replacement is the exception: the speaker withdraws a commitment already made and
puts another one in its place, so it is the only type that involves a semantic
update.

%% TODO(P3): the Switchboard composition figure (over half of annotated
%% disfluencies are simple repetitions, \citealp{charniakEditDetectionParsing2001,
%% shribergDISFLUENCIESSWITCHBOARD1996}) belongs with the data-availability
%% premise near RQ2, not here. Cite composition, not the 5.9%% rate.

The dominant computational treatment is detection-and-deletion, where the
reparandum is treated as noise and the goal is to recover a fluent string before
downstream tasks begin \citep{johnsonTAGbasedNoisychannelModel2004,
zayatsRobustCrossdomainDisfluency2018, rohanianBestBothWorlds2021}. The two
types that involve no update lose nothing under it: a repetition leaves its
content standing elsewhere in the string, and a restart never asserted any in the
first place. That fit is part of why the approach has held up so well.

Replacement is where deletion costs something. The material it discards was
committed to and then withdrawn, and both halves of that are part of what the
discourse records. Comprehenders behave accordingly, holding the reparandum in
context \citep{brennanHowListenersCompensate2001}, where it stays available as an
antecedent later in the same dialogue
\citep{purverComputationalModelsMiscommunication2018}.
In \textit{``My favourite dessert is cherry pie, actually, banana bread. They are
both good desserts,''} a system that receives only the cleaned string has no
second candidate for the plural pronoun ``they" to pick up.

\citet{senSemanticParsingDisfluent2021} make the assumption concrete. Every
disfluent utterance in their data is paired with the meaning representation in slot-value format of
its fluent counterpart, so the reparandum leaves no trace in the target at all.
A parser trained this way scores 0.81 exact match on fluent input and 0.02 on
disfluent versions of the same utterances. Synthetic augmentation recovers much
of that loss, though not evenly. Filled pauses end up largely handled, while
repairs stay below 0.3 even at the highest augmentation setting, on a disfluent
test set small enough that the figures are indicative rather than exact.

What that target discards has been the object of study elsewhere.
\citet{foxCrosslinguisticStudySelfrepair2010} code recycling, their term for
repetition, against replacement in English, German and Hebrew, and the two divide
the vocabulary between them. Recycling returns to a function word in roughly four
cases out of five, while replacement substitutes a content word between a third
and a half of the time. Replacement is the rarer of the two in all three
languages, and the one that carries lexical content when it occurs. Clinical work meets the same imbalance
from the other side. \citet{beierDisfluenciesAge2023} find no age effect on
disfluency when the count is aggregated; separating the types shows repeats
moving slightly with age while repairs do not, and the two predicted by different
variables. \citet{pirinenComprehensiveAnalysis2023} report significantly more
revisions and abandoned utterances in autistic young adults than in controls,
with no group difference for multisyllabic repetition. In both, the variation the
study is after sits in a subtype an aggregate disfluency rate averages away.
%% CUT(P2): the ASR-smoothing point (\citealp{romanaFluencyBank2024}) and the
%% dementia-detection cost (\citealp{farzanaHowYouSay2022}) were removed here.
%% Both concern the speech-to-text stage, which sits upstream of everything this
%% thesis does. Restore as a single subordinate clause if the scope argument is
%% wanted back.

Rarity of this kind sets how often a system meets the phenomenon, while the cost
of mishandling it is set by what a single instance carries. What a single
instance carries in both studies is a contrast: Fox et al. code the replaced item
against the item that replaced it, and a revision counts as a revision only
relative to what it revised. Neither half of the pair is the unit, so keeping the
reparandum available and keeping its relation to the repair are one requirement
rather than two.

Formal and computational semantics keep the retracted material. Standard DRT
\citep{Kamp1995} updates context monotonically, so nothing once added can be
taken back out, and taking a commitment back out is what a correction requires.
\citet{vanleusenIncompatibilityContextDiagnosis2004} supplies both a diagnosis
and a procedure, stated in Logical Description Grammar rather than in DRT itself:
a correction is felicitous when the corrective claim is incompatible with its
antecedent in context, and the update it triggers retracts that antecedent before
it adds the claim. Retraction is the operation on the context, and correction is
the discourse move that calls for it. Later work takes up the procedure and
leaves the formalism. Layered DRT demotes the offending layer
\citep{maierDenialCorrectionLayered2007}, and SDRT relates the two propositions
through a \textsc{Correction} rhetorical relation
\citep{lascaridesSegmentedDiscourseRepresentation2007}. What the three share is
the unit they work on. In each of them the material denied is a complete
assertion, and none reaches inside one. SDRT is the case that matters here, since
SBN takes its relation inventory from it. Applied to \textit{``I need a train, I
mean, a flight, to Berlin tonight,''} SDRT has to split one continuous utterance
into two speech acts before it can relate them. Whether the incompatibility
condition carries over to a speaker correcting themselves is a question
\citet{vanleusenIncompatibilityContextDiagnosis2004} leaves open.
Dynamic Syntax with TTR operates at the right granularity, marking reparandum
edges as repaired rather than deleting them
\citep{houghProcessingSelfrepairsIncremental2012, eshghiLearningGenerateCorr2023},
but maintaining incremental state is computationally expensive and does not fit
the sequence-to-sequence paradigm most current parsers are built on.

What is missing is a notation that can express non-monotonic self-repair at the
sub-propositional level, keep the reparandum structurally available, and still
be produced by a sequence model. Simplified Box Notation (SBN) \citep{bosSBN}
is a natural candidate for closing this gap. It is a linear serialisation of DRS,
inheriting DRT's box and scope semantics while remaining a string that a
sequence-to-sequence language model can produce directly. \citet{zhangEllipsisResolution2025} examine ellipsis in the same ecosystem with a
challenge set of 120 items, where PMB already supplies the target representation. %% TODO: add to .bib
Self-repair has no such representation, and building one has to come before any
diagnosis of parser behaviour.
%% TODO(P2): Sen \& Groves (2021) moves up into the detection-and-deletion paragraph,
%% carrying the 0.02 / 0.26 / 0.75 figures.

The thesis makes four contributions.

\begin{enumerate}[itemsep=0.2em]
  \item A representation of replacement self-repair in SBN, in which a
        \sbnterm{CORRECTION} box quarantines the reparandum and a sibling
        \sbnterm{CONJUNCTION} merges the repair into the context that was live
        before it. The scope of the treatment is stated explicitly, including the
        cases it cannot reach.
  \item The first instantiation of \sbnterm{CORRECTION}, a separator that SBN's
        inventory reserves but that appears nowhere in the annotated data of
        PMB 5.1.0.
  \item A pipeline that synthesises repair-aware (NL, SBN) pairs by transforming
        gold meaning representations directly, so that no parser output enters the
        training targets.
  \item A pilot study of how a parser trained on standard PMB data behaves on
        disfluent input, which shows the reparandum surviving into a well-formed
        graph alongside the repair.
\end{enumerate}

Three research questions follow.

\begin{enumerate}
  \item \textbf{Theory:} How should replacement self-repair be represented within SBN, and what is the coverage of that representation?
  \item \textbf{Data:} Given the scarcity of replacement self-repair in existing meaning-banked corpora, how can high-quality SBN-annotated data be synthesised from existing resources such as the Parallel Meaning Bank (PMB)?
  \item \textbf{Model:} Once fine-tuned on this synthesised data, how well do pretrained language models learn to parse replacement self-repair, and what error patterns remain?
\end{enumerate}

\section{Self-Repair Formalism in SBN}
\label{sec:formalism}

SBN represents a sentence as a sequence of boxes linked by separators, with
concepts and their argument roles nested inside each box. Some separators are
structural, such as \sbnterm{NEGATION} and \sbnterm{CONJUNCTION}. The rest are
discourse relations, which Bos takes over from SDRT more or less as they stand
\citep{bosSBN}. Two kinds of index hold the sequence together. Angle brackets on a
separator point between boxes, and signed numbers on a role point between
concepts. A signed index counts concepts along the whole sequence, and a box
boundary does not reset that count. A concept may only refer into a box that
dominates it, and where a role would otherwise cross that boundary the
\emph{-Of} form of the role turns the index backwards. That is what lets both
candidates in a repair attach to the verb in the box above them instead of
pointing forward across a box boundary. Inversion is part of the notation, though
the inverted forms this project needs, including \emph{ThemeOf} and
\emph{TimeOf}, do not occur in PMB's annotated data.

A correction does two things at once. In SDRT it denies one proposition and
commits to another, which can be written
\textsc{Correction}$(\pi_1,\pi_2) \equiv \neg K_{\pi_1} \wedge K_{\pi_2}$
\citep{lascaridesSegmentedDiscourseRepresentation2007}, and van Leusen's procedure
adds that the two halves update the context in different ways, one by retraction
and one by ordinary merge. \sbnterm{CORRECTION} is already available for the job,
since it entered SBN with the rest of SDRT's inventory, although PMB 5.1.0 never
uses it. One separator cannot carry the whole relation. A separator opens a single
box at the point where it stands, takes no arguments of its own, and has no way to
express a condition of two conjuncts that update differently. The two halves are
therefore split across two separators. \sbnterm{CORRECTION} opens the box that
quarantines the reparandum and carries the denial, while a sibling
\sbnterm{CONJUNCTION} merges the repair into the context that was live before the
correction opened and carries the commitment. Both boxes hang from the same
parent.

The same relation covers the cross-turn case without the second separator, and
the reason lies in what is being denied. When one speaker corrects another, both
arguments are complete contexts already, so a single separator relates two boxes
that exist anyway and no clause is interrupted. When a speaker corrects
themselves, the denied material is a constituent of a clause still in progress. It
has no box, so one has to be opened for it, and opening it breaks off the clause,
which then has to be resumed. The sibling \sbnterm{CONJUNCTION} is what that
resumption looks like. It appears exactly when the denied material sits below the
level of a proposition, which is also what tells the two configurations apart.

In \textit{``I ordered a banana bread, I mean, a cherry pie,''} both candidate
objects attach to the ordering event through an inverted \emph{ThemeOf} role,
so the reference points backward into \sbnterm{Box 0} instead of forward into
the boxes that quarantine and merge them:

\begin{lstlisting}[style=sbn]
person.n.01 EQU speaker              % I
time.n.08 TPR now
order.v.01 Agent -2 Time -1          % ordered
    CORRECTION <1
banana_bread.n.01 ThemeOf -1         % a banana bread,
    CONJUNCTION <2                   % I mean,
cherry_pie.n.01 ThemeOf -2           % a cherry pie
\end{lstlisting}

Both candidates reach the verb by counting backwards, which works because the
verb precedes them in the sentence and so also in the sequence. Where the
repaired concept comes first instead, most often a repaired subject, the
predicate that governs it is still ahead and no backward index reaches it. A
dummy concept, \sbnterm{entity.n.01}, is placed before the \sbnterm{CORRECTION}
to resolve this. The predicate's role lands on the dummy, which sits in the box
that dominates both candidates, and each candidate ties itself to the dummy with
\sbnterm{EQU}. This is the device Bos uses where a discourse referent and its
conceptual predicate end up in different boxes \citep{bosSBN}; here it gives the
two candidates a single referent the predicate can also see. A repaired subject
is among the commonest sites the notation has to cover, so this is not a
marginal use.
%% Source: Bos (2021), Variable-free DRS, Fig. 3 (p.5). Make sure the bib key
%% resolves to the 2021 paper and not the 2023 sequence-notation one.

The same dummy reaches repairs that target a role rather than a concept, such as
a repaired preposition, where it is introduced to carry the role and tied to the
event in the same way. That extends the device to material which has no concept
of its own.

In SBN the sub-propositional boundary falls at the separator. The mechanism
covers replacement below that level, as long as the reparandum and the repair
sit at the same level as each other.
Within that boundary, the notation handles repair on subjects, verbs, objects,
tense, aspect, negation scope, and adjuncts, and repair whose reparandum remains
available as an antecedent later in the discourse. The same boundary excludes
cases where the thing repaired is itself a separator. Replacing a quantifier would
mean undoing the nested \sbnterm{NEGATION} boxes that already encode it, and a
rhetorical relation or a modal verb runs into the same difficulty, since none of
the three is a concept and a \sbnterm{CORRECTION} box can quarantine nothing
else. The dummy concept works inside that limit rather than
around it, giving a repaired role somewhere to attach without making a separator
quarantinable. A further limitation is independent of the boundary.
\sbnterm{CORRECTION} can only quarantine material that follows it in the sequence,
so a repair whose reparandum was uttered several concepts earlier cannot be
captured without reordering.
%% TODO: an intra-turn repair followed by a cross-turn correction in the same
%% discourse was listed here as covered. Held back until KB 8.1 settles which box
%% the cross-turn CORRECTION attaches to.

\section{Pilot Study}
\label{sec:pilot}

Whether the notation needs that place is a question about what parsers do without
it. A model trained only on fluent data has nowhere to record a withdrawn
commitment, so a replacement in its input leaves it two moves: drop the
reparandum, which is what detection-and-deletion would do, or write it in as
though the speaker had asserted it. The pilot establishes which move a current
parser makes, and how much the choice costs.

The parser is Gemma-2-9B, fine-tuned with LoRA on PMB 5.1.0 train gold (NL, SBN)
pairs that contain no self-repair, and it was run over 836 sentences from
the same gold set under seven conditions: the sentence
unmodified (\textsc{gold}), a synthetic replacement inserted at the head,
middle, or tail of the sentence (\textsc{repair}), and the same three
positions with the interregnum \textit{``I mean''} inserted before the repair
(\textsc{repair+interr.}). These are the 865 sentences of nine words or more,
minus 29 items that turned out to be multi-sentence or direct speech; the length
threshold is what makes the three insertion positions distinguishable. The
sentences come from the training data by design. The gold condition shows what
the parser produces when the target for that sentence was among its training
examples, and the disfluent conditions are read against that. The main experiment
does not evaluate on training material, because generalisation is what it has to
measure. Every condition is scored against the same gold SBN graph for that sentence, so
the input string is the only thing that varies. Scoring uses smatch++
(citehere), which aligns the predicted and gold graphs and reports precision and
recall over the matched triples. The training steps are set to be aligned with Xiao
(2025), the source of the fine-tuning setup this pilot reuses.

\begin{table}[h]
\centering
\small
\begin{tabular}{lccc}
\toprule
condition & success \% & F1 (successes only) & F1 (penalized) \\
\midrule
gold (no repair) & 97.1 & 0.975 & 0.947 \\
repair & 92.0--95.5 & 0.825--0.842 & 0.760--0.804 \\
repair + interregnum & 70.5--86.6 & 0.779--0.805 & 0.557--0.697 \\
\bottomrule
\end{tabular}
\caption{Parser performance by condition (repair and repair+interregnum rows
range over head/mid/tail). Penalized F1 scores a failed parse as 0. Repair
variants were generated with a language model; the head and mid variants were
checked by hand, the tail variants were not.}
\label{tab:pilot-summary}
\end{table}

Performance degrades in two steps. Penalized F1 falls from 0.947 to
0.760--0.804 once a plain replacement is inserted, and falls further to
0.557--0.697 once an interregnum is added before the repair. Only a reparandum is
inserted in these conditions, so the gold graph stays the correct target
throughout. A parser that discarded the inserted material and returned the clean
graph would score on the disfluent conditions as well as on the gold one, and
that is what a detection-and-deletion system is built to produce. The pilot would
have counted it as success.

The outputs show something else. The reparandum is neither discarded nor marked,
but written into the graph as an additional concept, which leaves the result well
formed and wrong. Gold triples survive as a subset of the larger graph, so recall
holds while precision drops, and F1 settles in a range that reads as partial
competence. Relative indexing is what makes this cheap. Inserting a concept at
the front of the sequence moves every later position along by one, but the
distances between concepts do not change, so every index in the gold parse stays
correct and the whole of it survives around the inserted node. The parser's output for \textit{``Jerry, Tom threw a rock at Mary, but it
didn't hit her''} scores an F1 of 0.883 against the gold graph for the
unmodified sentence. It begins as follows, with the remaining second clause
left out since it does not bear on the reparandum:

\begin{lstlisting}[style=sbn]
male.n.02 Name "Jerry"
male.n.02 Name "Tom"
throw.v.01 Agent -1 Time +1 Theme +2 Destination +3
time.n.08 TPR now
rock.n.01
female.n.02 Name "Mary"
% "Jerry" (the reparandum) is written in as a free-standing node with no role
% attaching it to the rest of the graph. Everything after it is the gold parse
% unchanged: the absolute positions move along by one, but the indices are
% relative, so Agent -1 still reaches Tom and none of them has to be rewritten.
\end{lstlisting}

A companion ablation separates what the pretrained model already knows about
repair from what the narrow SBN fine-tuning changed. The base model runs with no
adapter, and the prompt asks it to mark the reparandum, the interregnum and the
repair in plain text. Here the split runs between cued and uncued repair.
Interregnum repair is detected at ceiling, 100\% across all
positions, because \textit{``I mean''} is a fixed lexical cue. Plain
replacement offers no such cue and is detected only 30.7--64.2\% of the time,
with the rate varying by where the reparandum sits. Detection therefore does not follow from a lexical cue alone.
It requires telling the reparandum from the repair by content, and even that is not the end of it. A
parser that has told them apart still has only two moves available in SBN as it
stands. It can drop the reparandum, which loses material the discourse may still
need, or write the reparandum in, which credits the speaker with a commitment they
withdrew. The outputs above take the second. There is no third move, because the
notation has nowhere to record that something was retracted.
Section~\ref{sec:formalism} supplies that place, and it has not been set out in
full or evaluated before.

\section{Methodology}

\subsection{Data}
\label{sec:data}

%% TODO: add to .bib -- passaliLARD2022, guptaDisflQA2021, abzianidzePMB2017,
%% fellbaumWordNet1998, schulerVerbNet2005, bakerFrameNet1998,
%% laurerLessAnnotating2024 (the NLI model the distinctness filter runs on).

Supplying that place is not enough, because a parser learns from examples and the
examples do not exist. PMB 5.1.0 pairs 9,552 English sentences with hand-checked
SBN graphs, and not one of them contains a self-repair. The corpora that do
contain self-repair carry no meaning representation, whether the disfluency is
natural, as in Switchboard, or synthetic, as in LARD and DISFL-QA. Where both
sides are present they do not meet: \citet{senSemanticParsingDisfluent2021} pair
every disfluent utterance with the representation of its fluent counterpart, so
the reparandum leaves no trace in the target. Nothing available pairs a
replacement with a representation that keeps both the reparandum and its relation
to the repair.

That pairing is built from PMB gold, and the construction is a single move. The
word the gold sentence already contains stays as the repair, since it records
what the speaker settled on, and a reparandum is manufactured in front of it. So
\textit{``The movie starts at ten o'clock''} becomes \textit{``The performance,
actually, movie starts at ten o'clock.''} The gold annotation survives inside the
result rather than being rewritten to accommodate it. Generation runs separately
over PMB's own train, dev and test partitions, and a source sentence yields at
most two variants. Every variant therefore stays in the partition its source came
from.

What a variant has to look like depends on what the repair does to the graph.
Substituting one content word leaves the rest of the graph standing, while
repairing a tense, an edge label or a name reaches parts of it that lexical
substitution never touches, and Section~\ref{sec:formalism} treats these as
separate cases. Table~\ref{tab:grid} is organised by that difference and not by
how long the repair runs on the surface. Most of its rows can be filled
automatically. Forwarding repair spans a subject and its verb at once, which the
transformation cannot express in one move, so the generator drafts those and a
person passes them. Repair whose reparandum is picked up later needs a second
sentence, and cross-turn correction needs a second speaker, so both are written
by hand.

\begin{table}[h]
\centering
\small
\begin{tabular}{llrrr}
\toprule
what the repair changes & how it is made & train & dev & test \\
\midrule
one content word                    & automatic          & 2,200 & 265 & 425 \\
a constant: name, number, time      & automatic          & 1,120 & 140 & 230 \\
tense or aspect                     & automatic          &   500 &  60 & 100 \\
an edge label                       & automatic          &   500 &  60 & 100 \\
a subject and its verb together     & drafted, reviewed  &   200 &  30 &  60 \\
reparandum reused later; cross-turn & hand-written       &   --- & --- &  80 \\
\midrule
total & & 4,520 & 555 & 995 \\
\bottomrule
\end{tabular}
\caption{Target composition, by the operation the repair requires of the graph.}
\label{tab:grid}
\end{table}

The repair corpus is added to PMB gold rather than replacing it, and comes to
about a third of the combined training data. That proportion answers the pilot: a
parser already inclined to write too much structure would only be pushed further
by a corpus in which repair is the norm. Supply is not what constrains it, since
the concept-substitution operator has produced some 15,200 candidates on the
training partition alone and the quotas sample from those. The hand-written rows
go to the test set by themselves, where they ask whether a model reaches a
construction it has never seen.

A site is available wherever the generator can find the word that realises a gold
concept, which it does by matching lemma and coarse part of speech against the
sentence. Where no match is found the site is passed over. The
\sbnterm{CORRECTION} box takes a single concept, so what gets withdrawn is a
single content word: \textit{the movie} against \textit{the performance}, never
\textit{the big dog} against \textit{the small cat}. Function words fall outside
that limit, since SBN gives them no concept of their own, and a reparandum may
repeat the determiner or the preposition in front of the content word and still
run to several words when spoken.

On the sentence, the reparandum goes in before the repair and both spans are
recorded. On the graph, \sbnterm{CORRECTION} opens the box that quarantines the
reparandum and a sibling \sbnterm{CONJUNCTION} merges the repair into the context
that was live before it, with both boxes hanging from the same parent. Because a
signed index counts concepts along the whole sequence, the insertion displaces
every reference reaching across it, and each of those is recomputed. In
\textit{p60/d2470} the object of \textit{put} is repaired, and the two candidates
attach to the verb through the inverted \emph{ThemeOf} of
Section~\ref{sec:formalism}:

\begin{lstlisting}[style=sbn]
% gold: put.v.01 Agent -1 Time +1 Theme +2 Destination +3
male.n.02 Name "Tom"
put.v.01 Agent -1 Time +1 Destination +4   % Theme is now inverted, below
time.n.08 TPR now
    CORRECTION <1
radar.n.01 ThemeOf -2                      % the attempt withdrawn
    CONJUNCTION <2                         % no wait,
thermometer.n.01 ThemeOf -3                % the gold concept, untouched
location.n.01 SZN +2
male.n.02 ANA -6                           % was -5
arm.n.01 PartOf -1
\end{lstlisting}
The displacement shows in the roles nobody edited: \sbnterm{Destination} moves
from \sbnterm{+3} to \sbnterm{+4} and \sbnterm{ANA} from \sbnterm{-5} to
\sbnterm{-6}, while what each of them points at is unchanged.

Whether the result is legal is settled before anything else is asked of it. Every
sample is read back with PMB's own parser and checked for four things: that the
graph has no cycle, that reparandum and repair sit in sibling boxes, that no edge
of the gold graph has been lost, and that no index has grown too long to be read.
Two of the four fail quietly, because an inserted box moves every index behind it
and arithmetic that is wrong still parses. Of the sites the generator attempts,
99.7\% come through legal. The surgery is not where the risk sits; the choice of
reparandum is.

A reparandum has to be a sibling of the repair in a contrast set, one slot with an
incompatible value in it, and what counts as the slot varies with the part of
speech. Ranking candidates by similarity is the obvious way to find one and it
returns the wrong relation, yielding near-synonyms at between a quarter and a
half of sites depending on the part of speech, where correcting \textit{watch}
to \textit{see} is pedantry rather than repair. The pool is assembled from WordNet co-hyponyms, VerbNet classes and
FrameNet frames, which encode contrast where distance encodes proximity.

Legality says nothing about whether a sentence could have been spoken.
\textit{``Adversary, Hitler assumed power in 1933''} passes every check above.
Whether a candidate differs from the repair enough to make a correction worth
uttering is a question about entailment, and natural language inference answers it
where embedding distance cannot: distance puts \textit{cheap} beside
\textit{expensive}, and an antonym is among the best reparanda available. Whether
the candidate could have stood in that position at all is a separate question, and
a language model's surprisal answers that one, compared among the candidates at a
single site rather than against a fixed threshold. A usable reparandum clears
both, being plausible where it was put without being interchangeable with what
replaced it.

Neither threshold can be argued for in advance, so both are being calibrated
against 150 pairs labelled by hand with the scores hidden. What comes out is a
band, since a candidate can fail by sitting too close to the repair as easily as
by sitting too far from it. The distinctness filter is implemented; the fit filter
and the band the two define together are still in calibration.

The interregnum varies on a second axis, and it is deliberately not spread evenly.
The pilot detected interregnum-marked repair at ceiling and missed unmarked repair
most of the time, so the unmarked case is the one a corpus has to teach and it is
generated in full. Five editing phrases take the place of the pilot's single
\textit{``I mean,''} which left open whether a model had recognised a repair or two
words. Phrases that signal reformulation are kept out, since \textit{``that is''}
and \textit{``or rather''} leave both versions standing and withdraw no
commitment. An interregnum changes the graph as well as the sentence, so a marked
sample and its unmarked twin are two records joined by an identifier.

Each record carries the identifier of the sentence it came from, the clean and
disfluent sentences, the clean and repair-aware graphs, the spans of reparandum,
repair and interregnum, and the operation the repair required. Reparandum and
repair are recorded rather than recovered, where the pilot had to find them by
diffing one graph against another. The outcome of each structural check is kept
with the record, so a defect can be inspected and not only filtered away.

Because the repair is spliced onto a gold graph by a transformation whose
parameters are known for every sample, the splice can be undone. Removing the
reparandum with its two boxes and running the renumbering backwards returns the
original graph without a parser, which lets one corpus be scored against a
repair-aware target, against the clean meaning, and against what a
detection-and-deletion system would be expected to produce. Failure can then be
localised, since a parse may be wrong about the repair or about the sentence
underneath it. The inverse transformation is not yet written.

Working now are the concept-substitution operator, the candidate pool, the graph
transformation and the structural checks, and the first calibration set is
labelled. Operators for constants, for tense and aspect and for edge labels are
still to be written, along with the inverse transformation and the fit threshold.
Generation at full scale follows those, then sampling to the quotas above, review
of the forwarding drafts, and the hand-written cases. Human work across the
pipeline comes to some 880 judgements, of which the calibration labels are 150 and
the rest is spot-checking the finished corpus, passing the drafted cases and
writing the test-only ones. Section~\ref{sec:timeline} places these in order.


\subsection{Model}
\label{sec:model}

Once that corpus exists, the question is what a parser trained on it does. The
pilot has already shown what one does without it. Given a replacement it had no
notation for, the model wrote the reparandum in as an ordinary concept and left
the graph well formed and wrong [hx: here I want to use a quick sentence to refresh the readers' memory that without trainig, the parser will sacrefise semantic accuracry to glue components together to make the output is structurally legal.]. 
Section~\ref{sec:formalism} supplies the missing
notation and Section~\ref{sec:data} the missing examples. What a parser does when
it has both is still open.

Only scale and training data are allowed to vary. Gemma-2 is released at 2B, 9B
and 27B, and each size is fine-tuned twice, once on PMB gold alone and once on
PMB gold with the repair corpus added. Sizes within one release share a
tokeniser, a pre-training recipe and an architecture, so a difference between
them is one of scale and not of provenance. The LoRA configuration and the
training schedule follow Xiao (2025), unchanged across all six runs. That keeps
the comparison between the two data conditions steady at every size. The 27B
model does not fit on a single 40 GB device in half precision. That pair of runs
therefore depends on a larger allocation, with a quantised base as the fallback.
%% TODO: resolve "Xiao (2025)" to a bib key; it is cited the same informal way
%% in Section~\ref{sec:pilot}.

Every output is checked for legality before it is scored, by the tests the
generator already runs on its own. Those tests reject a string that will not
parse, a graph carrying a cycle, a concept referring to itself, and an index
reaching past the end of the sequence. [hx: can we change it to an umbrealla term such as ill-formed ouputs] The pilot reported a success rate alongside two F1
figures, one over the successful parses and one counting failures as zero. The
same three columns are kept here.

Scoring runs against the repair-aware graph with smatch++, which aligns the
output to the target and reports precision and recall over the matched triples.
That figure stands as the headline result for each run. The same alignment
supports a finer reading. Every gold triple sits either inside the two boxes the
splice introduced or in the material around them, and the transformation recorded
which. Recall over the first group says whether the model found the reparandum
and the repair and put them in the right relation. Precision and recall over the
second say whether it got the sentence underneath right.

That second group is informative even where a model produces no
\sbnterm{CORRECTION} box at all, and the pilot outputs are the case in point.
Inserting a reparandum shifts every later concept along by one, but SBN indices
are relative, so the distances between concepts hold and the gold parse survives
around the inserted node. Recall stayed high there while precision fell.
Reading the two groups apart turns that observation into a measurement. It
separates a model that misses a repair from one that also damages the clause the
repair sits in.

Marked and unmarked variants of the same repair share a target graph, since an
editing phrase contributes no concept. The corpus joins the two by an identifier,
so their scores can be compared one record against its twin. A model that has
learned to pass over the editing phrase scores alike on both, and a gap between
them says the phrase is being read as content. The pilot gives reason to expect
one, since its parser lost ground at every insertion position once an interregnum
was added.

Those scores break down further by the operation each repair required, along the
categories of Table~\ref{tab:grid}. The breakdown locates trouble and is not
reported cell by cell. Several of those categories are small once the interregnum
split is applied, and a per-cell F1 over a few dozen items carries little weight.
Adding repair data to the training mix could also cost something on ordinary
sentences, so every repair-trained model is run on the unmodified PMB gold test
set as well. Its score there is read against the sibling trained on gold alone,
which isolates what the added data cost.

The hand-written items pass through the legality checks and the headline figure
in the same way. Grouping their gold triples needs a definition of its own for
each of the two kinds. A repair whose reparandum is picked up later has the
intra-sentential structure of every other item, so its repair-related triples are
the ones in the two spliced boxes. What it adds is a pronoun in the following
sentence that has to resolve to the reparandum, and that resolution is checked
separately. Cross-turn correction has a different shape, since both of its
arguments are complete contexts and a single \sbnterm{CORRECTION} relates them
with no sibling \sbnterm{CONJUNCTION}. The repair-related group is defined there
over the two boxes the relation connects. Neither kind appears in training at any
size.

\section{Timeline}
\label{sec:timeline}

%% NB: claims live in the contribution list in the Introduction.
%% This section lists concrete artefacts only (dataset, annotation manual,
%% adapter weights, evaluation toolkit, coverage report).
\section{Expected Outcomes And Deliverable}

\bibliography{}

\end{document}

